This work presents the design and implementation of a \textbf{Personal Behavior Anomaly Detection System (PBADS)} based on a microservices architecture. The system analyzes users' daily data (sleep, physical activity, diet, mood) and detects behavioral anomalies in real-time using unsupervised machine learning algorithms.

The solution is built on a \textbf{microservices architecture} with 13 distinct services, including an authentication service (JWT), a data service, an ML model service, an inference service for anomaly detection, an alert service, and a recommendation service powered by Google Gemini AI. Inter-service communication uses Kafka for asynchronous messaging and HTTP/REST for synchronous calls, with a fallback mechanism to ensure resilience.

The data model structures the main entities: \textit{User}, \textit{DailyData}, \textit{Alert} and \textit{Recommendation}. The system uses anomaly detection algorithms (Isolation Forest and One-Class SVM) to identify deviations from normal behaviors, generates alerts with severity levels (NORMAL, LOW, MEDIUM, HIGH), and provides personalized recommendations through integration with Google Gemini AI.

The observed benefits include: \textbf{automatic anomaly detection}, \textbf{personalized recommendations}, \textbf{real-time notifications}, and \textbf{interactive visualization} of behavioral data through a dashboard with temporal filters. The system offers high scalability through the microservices architecture and fault tolerance via Resilience4j.

\textbf{Keywords:} Anomaly Detection, Machine Learning, Microservices, Spring Boot, Kafka, PostgreSQL, Isolation Forest, One-Class SVM, Google Gemini AI, Personalized Recommendations, Microservices Architecture.
